Why do we take the complexity class $P$ as our main basis? There are three main reasons to do so, described as follows:
\begin{enumerate}
    \item \textbf{P is robust} \\ Polynomials seem to be the smallest class of bounds which makes $P$ robust.
    \item \textbf{Exponents are often small} \\ The exponent $c$ bounding the time of any Turing Machine capable to decide (solve) a language $L$ (decision problem) such that $L \in \textbf{P}$ can be huge. Anyway, for many problems of interest and in \textbf{P}, there are Turing Machines working with quadratic or cubic boundaries.
    \item \textbf{Nice closure properties} \\ It is relatively easy to prove that a given problem or task is \textbf{in} the class. Without constructing a Turing Machine explicitly, we can design an algorithm that guarantees the resolution of the problem in a polynomial running time.
\end{enumerate} 